% ============================================================
%  SECCIÓN — CÁLCULOS GEOTÉCNICOS (MURO)
% ============================================================
\section{Informe de cálculos geotécnicos}
\label{sec:muro_geotecnia}

% -----------------------------------------------------------
\subsection{Tensiones verticales efectivas por estrato}
\label{subsec:tensiones_verticales}
% -----------------------------------------------------------

El perfil estratigráfico del Sondeo~1 presenta cuatro horizontes diferenciados.
El nivel freático (N.F.) se localiza a $-2{,}50$~m de profundidad. Para la
obtención de las tensiones verticales efectivas $\sigma'_v$ se aplica el
\textbf{principio de Terzaghi}:

\begin{equation}
  \sigma'_v = \sigma_v - u
  \label{eq:terzaghi}
\end{equation}

donde $u = \gamma_w \cdot h_w$ es la presión intersticial, con
$\gamma_w = 1{,}00$~Tn/m³ y $h_w$ la altura de agua sobre el punto considerado.
Por encima del N.F.\ la presión intersticial es nula ($u=0$), por lo que
$\sigma'_v = \sigma_v = \gamma \cdot z$.

\begin{supuestobox}[Densidades adoptadas]
  \begin{itemize}
    \item Estrato~1 no saturado ($0 \to -2{,}50$~m):
          $\gamma_{\text{ap}} = 1{,}82$~Tn/m³
    \item Estrato~1 saturado ($-2{,}50 \to -3{,}50$~m):
          $\gamma_{\text{sat}} = 1{,}98$~Tn/m³~$\Rightarrow$~%
          $\gamma_{\text{sum}} = \gamma_{\text{sat}} - \gamma_w
           = 1{,}98 - 1{,}00 = 0{,}98$~Tn/m³
    \item Estrato~2 ($-3{,}50 \to -8{,}00$~m):
          $\gamma_{\text{sat}} = 1{,}78$~Tn/m³~$\Rightarrow$~%
          $\gamma_{\text{sum}} = 0{,}78$~Tn/m³
    \item Estrato~3 ($-8{,}00 \to -25{,}00$~m):
          $\gamma_{\text{sat}} = 1{,}92$~Tn/m³~$\Rightarrow$~%
          $\gamma_{\text{sum}} = 0{,}92$~Tn/m³
  \end{itemize}
\end{supuestobox}

\subsubsection{Cálculo punto a punto}

\paragraph{Cota 0,00~m (superficie):}
\begin{equation*}
  \sigma_v = 0{,}00~\text{Tn/m}^2 \qquad u = 0{,}00~\text{Tn/m}^2
  \qquad \resultado{\sigma'_v = 0{,}00~\text{Tn/m}^2}
\end{equation*}

\paragraph{Cota $-2{,}50$~m (techo N.F., base estrato~1a):}
\begin{align*}
  \sigma_v   &= \gamma_{\text{ap}} \cdot h_1
               = 1{,}82 \times 2{,}50 = 4{,}55~\text{Tn/m}^2 \\
  u          &= 0{,}00~\text{Tn/m}^2 \\
  \sigma'_v  &= \resultado{4{,}55~\text{Tn/m}^2}
vale \end{align*}

\paragraph{Cota $-3{,}50$~m (base estrato~1b / techo estrato~2):}
\begin{align*}
  \sigma_v   &= 4{,}55 + \gamma_{\text{sat},1b} \cdot h_{1b}
               = 4{,}55 + 1{,}98 \times 1{,}00 = 6{,}53~\text{Tn/m}^2 \\
  u          &= \gamma_w \cdot 1{,}00 = 1{,}00~\text{Tn/m}^2 \\
  \sigma'_v  &= 6{,}53 - 1{,}00 = \resultado{5{,}53~\text{Tn/m}^2}
\end{align*}

\paragraph{Cota $-8{,}00$~m (base estrato~2 / techo estrato~3):}
\begin{align*}
  \sigma_v   &= 6{,}53 + \gamma_{\text{sat},2} \cdot h_2
               = 6{,}53 + 1{,}78 \times 4{,}50 = 6{,}53 + 8{,}01 = 14{,}54~\text{Tn/m}^2 \\
  u          &= \gamma_w \cdot 5{,}50 = 5{,}50~\text{Tn/m}^2 \\
  \sigma'_v  &= 14{,}54 - 5{,}50 = \resultado{9{,}04~\text{Tn/m}^2}
\end{align*}

\paragraph{Verificación por densidad sumergida (coherencia):}
\[
  \sigma'_v(-8{,}00) = \sigma'_v(-3{,}50)
    + \gamma_{\text{sum},1b} \cdot 1{,}00
    + \gamma_{\text{sum},2} \cdot 4{,}50
  = 4{,}55 + 0{,}98 + 3{,}51 = 9{,}04~\text{Tn/m}^2 \quad\checkmark
\]

\begin{table}[H]
  \centering
  \caption{Resumen de tensiones verticales efectivas — Sondeo~1}
  \label{tab:tensiones_sondeo}
  \begin{tabular}{@{}lccccc@{}}
    \toprule
    \textbf{Cota (m)} & \textbf{Estrato}
    & $\sigma_v$ \textbf{(Tn/m²)}
    & $u$ \textbf{(Tn/m²)}
    & $\sigma'_v$ \textbf{(Tn/m²)} \\
    \midrule
    0{,}00           & —           &  0{,}00 & 0{,}00 & 0{,}00 \\
    $-$2{,}50        & Arena 1a    &  4{,}55 & 0{,}00 & 4{,}55 \\
    $-$3{,}50        & Arena 1b    &  6{,}53 & 1{,}00 & 5{,}53 \\
    $-$8{,}00        & Arcilla 2   & 14{,}54 & 5{,}50 & 9{,}04 \\
    $-$25{,}00       & Arcilla 3   & 47{,}30 &22{,}50 &24{,}80 \\
    \bottomrule
  \end{tabular}
\end{table}

% -----------------------------------------------------------
% Los coeficientes de Rankine y el cálculo detallado de empujes activos
% se desarrollan en la sección~\ref{subsec:empujes_rankine} del capítulo
% anterior, donde se obtiene E_T = 25,78 Tn/ml y M_T = 74,77 Tn·m/ml.
% -----------------------------------------------------------
